\section{Primero, definir el problema: qué significa aquí «razonamiento»}

El debate sobre si «los modelos de lenguaje grandes realmente razonan» suele mezclar filosofía, ciencia cognitiva, lógica formal y capacidades de producto. Al inicio, Denny Zhou acota el problema de forma deliberada: esta clase estudia solo un objeto operable y medible, la \textbf{secuencia de token intermedios} que el modelo genera entre la entrada y la respuesta final. Esta definición no exige que la trayectoria se parezca al cerebro humano ni garantiza que explique con fidelidad el cómputo interno del modelo; solo pregunta si agregar pasos de generación intermedios visibles amplía el cómputo que el modelo puede realizar, modifica la señal de entrenamiento y mejora la selección de la respuesta final.

\lecturefigure{slide-01.jpg}{Título del curso: LLM Reasoning}{1}

\teachervoice{00:00:52--00:02:30, el docente afirma explícitamente que no participa en el debate sin definición sobre si «los LLM realmente razonan o no». Esta clase fija reasoning como los intermediate tokens entre input y output; toda la teoría, la decodificación y los métodos de entrenamiento posteriores se desarrollan en torno a esta definición operativa.}

\subsection{Los token intermedios son una interfaz de cómputo, no una prueba de personalidad}

Sea la entrada $x$, la secuencia de razonamiento intermedio $r=(r_1,\ldots,r_T)$ y la respuesta final $a=(a_1,\ldots,a_M)$. Un modelo autorregresivo asigna una probabilidad a toda la respuesta:
\[
p_\theta(r,a\mid x)
=\prod_{t=1}^{T}p_\theta(r_t\mid x,r_{<t})
 \prod_{j=1}^{M}p_\theta(a_j\mid x,r,a_{<j}).
\]
Aquí $\theta$ son los parámetros del modelo, $T$ es el número de token intermedios, $M$ es el número de token de la respuesta, y $r_{<t}$ y $a_{<j}$ denotan los prefijos ya generados. El cambio esencial no es que «el modelo de pronto tenga un monólogo interior», sino que la generación de la respuesta, que originalmente ocurría en una sola pasada hacia adelante, se despliega en $T+M$ pasos de cómputo condicional; cada token nuevo puede leer los resultados anteriores y escribir un estado temporal de vuelta en el contexto.

\lecturefigure{slide-02.jpg}{Definición operativa de esta clase: los pasos intermedios entre la entrada y la salida}{2}

\begin{importantbox}{Definición operativa}
En esta clase, LLM reasoning se refiere a los token intermedios que el modelo genera antes de la respuesta final. Es una \textbf{trayectoria de cómputo observable}: puede usarse para continuar el cómputo, entrenar, filtrar y agregar; pero «legible» no equivale a «fiel», y «parecido a lo humano» no equivale a «mismo mecanismo».
\end{importantbox}

\begin{warningbox}{La sustitución conceptual más frecuente}
Al ver pasos en lenguaje natural, es fácil sustituir «la trayectoria ayuda a obtener la respuesta correcta» por «la trayectoria revela por completo el proceso causal interno del modelo». Esta clase respalda lo primero, pero de ello no se sigue automáticamente lo segundo. Para evaluar una trayectoria hay que distinguir, como mínimo, la corrección de la respuesta final, la validez de los pasos, la posibilidad de verificarlos y la fidelidad de la explicación.
\end{warningbox}

\subsection{La última letra: por qué diseñar tareas que no se resuelvan por atajos}

La concatenación de las últimas letras de «artificial intelligence» es «le». La respuesta directa tiene solo dos caracteres y no muestra cómo la obtuvo el modelo; la versión paso a paso primero extrae la última letra de artificial, $l$, luego la última letra de intelligence, $e$, y al final las concatena. El valor didáctico de esta pequeña tarea no está en su dificultad, sino en que descompone un proceso secuencial en estados temporales explícitos, lo que permite a los investigadores distinguir entre «haber memorizado un patrón frecuente» y «haber ejecutado la transformación paso a paso».

\lecturefigure{slide-03.jpg}{Concatenación de últimas letras: comparación entre la respuesta directa y los pasos intermedios}{3}

\teachervoice{00:03:18--00:05:10, el docente explica que al principio probaron con la «concatenación de primeras letras», pero la gran cantidad de siglas en la web permitía que incluso los modelos antiguos respondieran correctamente sin dificultad, por lo que cambiaron a las últimas letras. Lección práctica: en una tarea sintética hay que revisar los atajos presentes en el corpus de entrenamiento; de lo contrario, el benchmark podría estar midiendo memoria de coocurrencias y no el algoritmo objetivo.}

El trabajo de rationale generation de 2017 utilizó pasos intermedios en lenguaje natural para problemas algebraicos de enunciado; investigaciones neurosimbólicas anteriores solían expresar el proceso intermedio mediante programas, fórmulas lógicas o estados de búsqueda. El giro histórico no consiste en que «por primera vez hubiera estados intermedios», sino en colocar esos estados en la misma secuencia de lenguaje que la entrada y la salida, de modo que puedan generarse y aprenderse directamente con modelos de secuencias.

\lecturefigure{slide-04.jpg}{Rationale generation: el lenguaje natural como soporte del cómputo intermedio}{4}

\begin{knowledgebox}{Términos clave: rationale, scratchpad y Chain-of-Thought}
\textbf{rationale} pone el acento en la explicación o el fundamento de la solución; \textbf{scratchpad}, en usar un espacio intermedio escribible para realizar el cómputo; \textbf{Chain-of-Thought (CoT)} suele referirse a pasos consecutivos expresados en lenguaje. En la práctica, los tres pueden usar secuencias de token similares, pero sus objetivos de investigación difieren: la explicación, el espacio de trabajo para el cómputo y la capacidad de razonamiento no deben confundirse solo por su nombre.
\end{knowledgebox}

\subsection{Por qué los token intermedios aumentan la computabilidad}

Las diapositivas oficiales plantean una línea teórica: si un problema puede resolverse con un Boolean circuit de tamaño $T$, entonces un Transformer de tamaño fijo puede simular ese cómputo paso a paso generando $O(T)$ token intermedios; si se obliga al modelo a producir directamente la respuesta final, podría requerirse una red muy profunda, o incluso podría ser imposible expresarlo con la profundidad dada. Intuitivamente, la profundidad de la red aporta «capas secuenciales dentro de los parámetros» y la longitud de generación aporta «pasos secuenciales durante la inferencia».

\lecturefigure{slide-05.jpg}{Los token intermedios trasladan el cómputo secuencial de la profundidad de la red a la longitud de generación}{5}

El presupuesto total de cómputo puede escribirse de forma aproximada como
\[
C_{\text{total}}\approx C_{\text{prefill}}+T\,C_{\text{decode-step}}.
\]
Aquí $C_{\text{prefill}}$ es el costo de procesar la entrada, $C_{\text{decode-step}}$ es el costo marginal de generar un token nuevo y $T$ es la longitud de los pasos intermedios. Una trayectoria más larga aporta más cómputo secuencial, pero también aumenta la latencia, el uso de KV cache y la acumulación de errores; por eso, «permitir más token» no es una mejora gratuita, sino que convierte un problema de capacidad en un problema de presupuesto y de selección.

\begin{importantbox}{Conclusión central desde la perspectiva del cómputo}
El valor de los token intermedios consiste, ante todo, en \textbf{aumentar la profundidad de cómputo secuencial disponible durante la inferencia}. Esto explica por qué un modelo con los mismos parámetros puede resolver, cuando se le permite desplegar pasos, problemas en los que falla al responder directamente; también explica por qué test-time compute debe precisar «qué se está ampliando exactamente: la longitud, el número de muestras, la amplitud de búsqueda o las llamadas a herramientas».
\end{importantbox}

\subsection{Resumen de la sección}

Esta sección fija reasoning como token intermedios y explica su función mediante la descomposición probabilística, las tareas sintéticas y la profundidad de cómputo. La definición evita deliberadamente las analogías con la mente humana: lo que realmente queda por resolver es cómo encontrar buenas trayectorias, cómo entrenar su distribución, cómo agregar varias trayectorias y cómo demostrar que la respuesta final es confiable.
